\subsubsection{万能欧几里得}
把 linear 看成 U（穿过横轴）和 R（穿过纵轴）的操作序列，复合操作来优化复杂度。
\begin{lstlisting}
//Operation.
//Cross a horizontal line => U operation (impact of floor sum+1)
//Cross a vertical line => R operation (impact of i+1)
struct node{
    ll s,numU,numR;
    node(ll s=0,ll numU=0,ll numR=0){
        this->s=s,this->numU=numU,this->numR=numR;
    }
    node operator + (const node& b){
        node c;
        c.numR=numR+b.numR;
        c.numU=numU+b.numU;
        c.s=s+b.s+numU*b.numR;
        return c;
    }
};
//pow of operation
node opow(node a,ll b){
    node res;
    while(b){
        if(b&1)res=res+a;
        b>>=1,a=a+a;
    }return res;
}
// need a>=0, b>=0, b<m
// 计算 [1,n] 的UR复合操作的结果
node work(ll a,ll b,ll n,ll m,node U,node R){
    if(!n)return node();//构造函数应该取单位元，若没有单位元则需要小心处理opow里b=0和本行的情况。
    if(a>=m)return work(a%m,b,n,m,U,opow(U,a/m)+R);
    ll cu=(n*a+b)/m;//Warning of overflow, use long double if needed.
    if(!cu)return opow(R,n);
    ll cnt=n-(m*cu-b-1)/a;//Warning of overflow, use long double if needed.
    return opow(R,(m-b-1)/a)+U+work(m,(m-b-1)%a,cu-1,a,R,U)+opow(R,cnt);
}
node U,R;

// 使用例：\sum_{i=0}^{n-1} \lfloor\frac{ai+b}{m}\rfloor
// for b>=m, use opow(U,b/m)+work(a,b%m,n,m,U,R) to obtain the result. 
void solve(){
	U=(node){0,1,0},R=(node){0,0,1};
    n=d-1,m=d;ll a=d,b=d;
    node res=opow(U,b/m)+work(a,b%m,n,m,U,R);//b>=m的时候要线提出几个 U 来用。
    printf("%lld\n",res.s+b/m);
    // out(res.s+b/m);puts("");
}
\end{lstlisting}

